\subsection{Clarke dual action principle}\label{subsec:preliminaries-clarke-dual-action-principle}

The finite Haim--Kislev formula used in
Section~\ref{sec:quadratic-program-algorithm-hk2019} rests on a dual
variational viewpoint. The literature often states this connection in a fixed
\([0,1]\)-period normalization. We use a free-period normalization instead:
the dual curves below have period \(T\), satisfy \(A(z)=T\), and are not
centered by a condition such as \(\int z=0\). Translation is handled as a
symmetry, not as a constraint. This is equivalent to the normalized formulation
used by Haim--Kislev, but the numerical equalities below are written in the
free-period convention used throughout this thesis.

\begin{lemma}[Fenchel duality for the gauge]\label{lem:preliminaries-fenchel-duality}
The Legendre--Fenchel conjugate of \(g_K^2\) is
\(\frac{1}{4}h_K^2\). Equivalently,
\[
  g_K^2(x) + \frac{1}{4}h_K^2(y) \geq \langle x,y\rangle
\]
for all \(x,y\in\mathbb{R}^4\), with equality if and only if
\[
  y\in \partial g_K^2(x)
  \quad\Longleftrightarrow\quad
  x\in \partial\bigl(\tfrac{1}{4}h_K^2\bigr)(y).
\]
\end{lemma}

\begin{proof}
For fixed \(y\), write \(x=\lambda u\) with \(\lambda\geq0\) and
\(g_K(u)=1\). Since \(g_K(\lambda u)^2=\lambda^2\),
\[
  (g_K^2)^*(y)
  = \sup_{\lambda\geq0}
      \sup_{g_K(u)=1}\{\lambda\langle u,y\rangle-\lambda^2\}
  = \sup_{\lambda\geq0}\{\lambda h_K(y)-\lambda^2\}
  = \frac{1}{4}h_K(y)^2.
\]
The Fenchel inequality and the equality/subdifferential condition are the
standard equality cases for a convex function and its conjugate.
\end{proof}

Define the dual functional on absolutely continuous closed curves
\[
  z\colon[0,T]\to\mathbb{R}^4
\]
by
\begin{equation}\label{eq:preliminaries-dual-functional}
  I_K(z) \coloneqq
  \frac{1}{4}\int_0^T h_K^2(-J_0\dot{z}(t))\,dt.
\end{equation}
The dual problem is
\begin{equation}\label{eq:preliminaries-dual-problem}
  \min\{I_K(z):
    z\in W^{1,2}([0,T],\mathbb{R}^4),\;
    z(0)=z(T),\; T>0,\; A(z)=T\}.
\end{equation}

\begin{definition}[Weak dual critical curve]\label{def:preliminaries-weak-dual-critical}
A closed curve
\[
  z\in W^{1,2}([0,T],\mathbb{R}^4),
  \qquad A(z)=T,
\]
is \emph{weak dual critical} if there are
\(\nu\in\mathbb{R}\) and \(b\in\mathbb{R}^4\) such that
\begin{equation}\label{eq:preliminaries-dual-euler-lagrange}
  \nu z(t)+b
  \in \partial\bigl(\tfrac{1}{4}h_K^2\bigr)(-J_0\dot{z}(t))
  \quad\text{for a.e. }t.
\end{equation}
Equivalently,
\[
  -J_0\dot{z}(t)\in\partial g_K^2(\nu z(t)+b)
  \quad\text{for a.e. }t.
\]
\end{definition}

We use one general result from Clarke's nonsmooth variational theory:
minimizers of the dual action functional exist, and the nonsmooth Lagrange
multiplier rule applies to the action constraint \cite[Ch.~5]{Clarke1983}.
For nonsmooth convex bodies this is the variational input used by
Artstein--Avidan--Ostrover and Haim--Kislev
\cite[Propositions~2.5, 2.7 and Lemmas~5.1, 5.2]{AAO2014}
\cite[Section~2]{HK2017}.

Those sources state the dual problem in a fixed-\([0,1]\) normalization, with
closed loops \(w\) satisfying \(A(w)=\frac12\). The passage to the free-period
normalization above is by explicit rescaling. If \(z\) is defined on \([0,T]\)
and satisfies \(A(z)=T\), then
\[
  w(s)=\frac{1}{\sqrt{2T}}\,z(Ts)
\]
satisfies \(A(w)=\frac12\) and \(I_K(w)=I_K(z)/2\). Conversely, if \(w\) is a
fixed-period dual curve and \(c=2I_K(w)\), then
\[
  z(t)=\sqrt{2c}\,w(t/c),
  \qquad 0\leq t\leq c,
\]
satisfies \(A(z)=c\) and \(I_K(z)=c\). Under the same rescaling, the fixed-period
weak-critical equation used by Haim--Kislev becomes
\eqref{eq:preliminaries-dual-euler-lagrange}. If a cited version fixes the
translation freedom by imposing a centering condition on \(w\), this is only a
choice of representative: \(A\), \(I_K\), and the weak-critical equation are
translation-invariant after adjusting the constant \(b\). Thus the cited
analytic input is the fixed-period existence and weak-criticality theorem; the
rest of this subsection writes out the convention translation and Reeb
correspondence in the free-period, uncentered notation used here.

\begin{lemma}[Euler--Lagrange equation in thesis convention]
\label{lem:preliminaries-dual-euler-lagrange}
Every minimizer of the dual problem is weak dual critical in the sense of
Definition~\ref{def:preliminaries-weak-dual-critical}.
\end{lemma}

\begin{proof}
Let \(z\) be a minimizer of the free-period dual problem, and rescale it to the
fixed-\([0,1]\) curve \(w\) above. The fixed-period Clarke theorem gives the
nonsmooth Euler--Lagrange inclusion for \(w\), with one multiplier for the
constraint \(A(w)=\frac12\). Rescaling this inclusion back to \([0,T]\) gives
the same equation for \(z\). We spell out the resulting equation in the
free-period notation.

There is \(\alpha\in\mathbb{R}\) such that \(z\) satisfies the nonsmooth
Euler--Lagrange inclusion for
\[
  L(z,\dot z)
  =\frac14h_K^2(-J_0\dot z)
    +\alpha\langle -J_0\dot z,z\rangle.
\]
Thus there is an absolutely continuous \(\xi\) with
\[
  \xi(t)\in\partial_{\dot z}L(z(t),\dot z(t)),
  \qquad
  \dot\xi(t)=\nabla_z L(z(t),\dot z(t)).
\]
The right-hand side is \(-\alpha J_0\dot z\). On the left-hand side, the chain
rule for \(\frac14h_K^2(-J_0\dot z)\) gives
\[
  \partial_{\dot z}L
  =J_0\,\partial\bigl(\tfrac14h_K^2\bigr)(-J_0\dot z)
    +\alpha J_0z.
\]
Hence there is
\(p(t)\in\partial(\frac14h_K^2)(-J_0\dot z(t))\) with
\[
  J_0\dot p+\alpha J_0\dot z=-\alpha J_0\dot z.
\]
Since \(J_0\) is invertible, \(\dot p=-2\alpha\dot z\), and therefore
\[
  p(t)=b-2\alpha z(t)
\]
for some constant \(b\in\mathbb{R}^4\). With \(\nu=-2\alpha\), this is exactly
\eqref{eq:preliminaries-dual-euler-lagrange}. The equivalent formulation
follows from Lemma~\ref{lem:preliminaries-fenchel-duality}.
\end{proof}

\begin{lemma}[Scaling and translation]\label{lem:preliminaries-dual-symmetry}
Let \(z\) be weak dual critical on \([0,T]\). For every
\(\alpha>0\) and every \(b\in\mathbb{R}^4\), the curve
\[
  z'(t)=\alpha z(t/\alpha^2)+b,
  \qquad 0\leq t\leq \alpha^2T,
\]
is again weak dual critical. Moreover
\[
  A(z')=\alpha^2A(z)=\alpha^2T,
  \qquad
  I_K(z')=I_K(z).
\]
\end{lemma}

\begin{proof}
Closure is immediate. The action identity follows by substituting
\(s=t/\alpha^2\); the term containing the constant vector \(b\) integrates to
zero because \(z\) is closed. For the dual functional,
\[
  I_K(z')
  =\frac{1}{4}\int_0^{\alpha^2T}
    h_K^2\!\left(-\frac{1}{\alpha}J_0\dot{z}(t/\alpha^2)\right)\,dt
  =I_K(z),
\]
using the \(1\)-homogeneity of \(h_K\). The Euler--Lagrange inclusion is
preserved because
\(\partial(\frac14 h_K^2)(\lambda y)=\lambda
\partial(\frac14 h_K^2)(y)\) for \(\lambda>0\).
\end{proof}

\begin{theorem}[Clarke dual action principle in free-period form]\label{thm:preliminaries-clarke-dual-action}
For a smooth convex body \(K\), and for polytopes after interpreting closed
characteristics as generalized Reeb orbits in the sense of
Section~\ref{sec:generalized-reeb-orbits-polytopes}, the primal minimum and the
dual minimum \eqref{eq:preliminaries-dual-problem} agree:
\[
  c_{\mathrm{EHZ}}(K)
  = \min_\gamma A(\gamma)
  = \min_z I_K(z).
\]
More precisely, every contact-normalized closed characteristic is weak dual
critical, and every weak dual critical curve has a translate/rescale in the
sense of Lemma~\ref{lem:preliminaries-dual-symmetry} that is a
contact-normalized closed characteristic.
\end{theorem}

\begin{proof}
Let \(\gamma\) be a contact-normalized closed characteristic of period \(T\).
Then \(g_K^2(\gamma)=1\),
\(-J_0\dot{\gamma}\in\partial g_K^2(\gamma)\), and \(A(\gamma)=T\). By
Lemma~\ref{lem:preliminaries-fenchel-duality},
\[
  \gamma(t)\in
  \partial\bigl(\tfrac{1}{4}h_K^2\bigr)(-J_0\dot{\gamma}(t))
  \quad\text{for a.e. }t.
\]
This is the dual Euler--Lagrange condition with \(\nu=1\) and \(b=0\).
Thus \(\gamma\) is weak dual critical, and
Lemma~\ref{lem:preliminaries-dual-symmetry} gives the whole
translate/rescale family.

Conversely, let \(z\) be weak dual critical of period \(T\). Choose
\(\nu,b\) as in Definition~\ref{def:preliminaries-weak-dual-critical}. The
equivalent inclusion in that definition gives
\[
  -J_0\dot{z}(t)\in\partial g_K^2(\nu z(t)+b).
\]
The multiplier \(\nu\) is nonzero. If \(\nu=0\), then closedness of \(z\) and
convexity would imply \(0\in\partial g_K^2(b)\), hence \(b=0\), and then the
dual Euler--Lagrange condition would force \(\dot{z}=0\), contradicting
\(A(z)=T>0\).

Set \(x(t)=\nu z(t)+b\). For almost every \(t\), the chain rule for the convex
function \(g_K^2\) gives
\[
  \frac{d}{dt}g_K^2(x(t))
  = \langle -J_0\dot{z}(t),\nu\dot{z}(t)\rangle
  =0.
\]
Since \(z\) is closed, \(g_K^2(x(t))\) is constant; write
\(g_K(x(t))=c>0\). Integrating the Fenchel equality gives
\[
  c^2T+I_K(z)
  =\int_0^T\langle x(t),-J_0\dot{z}(t)\rangle\,dt
  =2\nu A(z)=2\nu T.
\]
The left-hand side is positive, so \(\nu>0\).

Define
\[
  \gamma(t)=\frac{\nu z(t/\nu)+b}{c},
  \qquad 0\leq t\leq \nu T.
\]
Then \(g_K(\gamma)=1\). Using the \(2\)-homogeneity of \(g_K^2\),
\(\partial g_K^2(x/c)=c^{-1}\partial g_K^2(x)\), so
\[
  -J_0\dot{\gamma}(t)
  =-\frac{1}{c}J_0\dot{z}(t/\nu)
  \in \partial g_K^2(\gamma(t)).
\]
Thus \(\gamma\) is a contact-normalized closed characteristic. A direct action
computation gives
\[
  A(\gamma)=\frac{\nu^2}{c^2}A(z)=\frac{\nu^2}{c^2}T.
\]
Since \(\gamma\) is contact-normalized, \(A(\gamma)=\nu T\). Hence
\(\nu=c^2\), and
\[
  \gamma(t)=c\,z(t/c^2)+\frac{b}{c},
\]
so \(\gamma\) lies in the translate/rescale family of \(z\).

It remains to compare the values. For any contact-normalized closed
characteristic \(\gamma\), Fenchel equality gives
\[
  g_K^2(\gamma)+\frac{1}{4}h_K^2(-J_0\dot{\gamma})
  =\langle \gamma,-J_0\dot{\gamma}\rangle.
\]
After integration,
\[
  T+I_K(\gamma)=2A(\gamma)=2T.
\]
Thus \(I_K(\gamma)=T=A(\gamma)\). The value \(I_K\) is constant on the
translate/rescale family by Lemma~\ref{lem:preliminaries-dual-symmetry}.
Conversely, the argument above sends every weak dual critical curve to a
contact-normalized closed characteristic in its translate/rescale family.
Therefore the infima over contact-normalized closed characteristics and over
weak dual critical curves agree. The nonsmooth existence part of Clarke's
principle gives a minimizer of the dual problem, and
Lemma~\ref{lem:preliminaries-dual-euler-lagrange} shows that such a minimizer
is weak dual critical. Hence the dual minimum is attained among the curves just
considered, and the primal and dual minima agree.
\end{proof}

The point of the theorem is not just the equality of two numbers. The primal
formulation constrains a curve to lie on \(\partial K\) and to satisfy a
differential inclusion. The dual formulation removes the position constraint
and makes the functional depend only on \(\dot{z}\). In their equivalent
fixed-period normalization, Haim--Kislev exploit this freedom to choose a
capacity minimizer on a polytope with a finite, simple facet-word description
\cite{HK2017}.
